\section{Horizon growth and the allocation of construction work}\label{sec:resources86}
A recursive endpoint query can expand into a descendant tree even when no state lattice is stored. The inherited vector experiment documents this effect. Compiling an attained continuation changes the dependency graph: a date-$t$ query evaluates the stored date-$(t+1)$ function rather than reopening its construction. This section quantifies that alternative while retaining its state, action and innovation costs.

At date $t$, let $N_t$ be the number of state--context cover nodes, $A_t$ the maximum number of feasible action representatives, $Q_t$ the number of integration representatives, and $C_{t+1}$ the arithmetic cost of a compiled continuation query. Let $P_t$ cover primitive evaluation, $R_t$ feasible repair and $D_t$ certificate arithmetic. An explicit account is
\begin{equation}\label{eq:work86}
 W\le\sum_{t<T}\bigl\{N_tA_tQ_t(C_{t+1}+P_t)
             +N_t(R_t+D_t)+\operatorname{Compile}_t\bigr\}
       +W_{\rm load}+W_{\rm output}.
\end{equation}
Memory for two rolling value layers is separate from the memory required to retain a complete policy. With tensor interpolation, $C_t=O(k2^k)$ after cell location; dense envelope queries have $C_t=O(kN_t)$. Exact rational work additionally depends on operand bit lengths. The displayed inequality is an operation account, not a conversion into seconds or energy.

\begin{proposition}[Compiled horizon bound]\label{prop:horizon86}
Suppose the primitive moduli and the effective cover, feasible-net and integration requirements are uniformly bounded in date. Suppose also $0<\beta_t\le\bar\beta<1$ and that the sum of the construction and acquisition allowances at each nonterminal date in \eqref{eq:constructionloss49} and \eqref{eq:acquire49}, as well as the terminal allowance, is at most $\eta$. Then the complete policy loss is at most $\eta/(1-\bar\beta)$, uniformly in the horizon. Choosing $\eta=(1-\bar\beta)\varepsilon$ permits date-uniform resolutions. If the terms in braces in \eqref{eq:work86} are at most $K(\varepsilon)$, construction work is at most $TK(\varepsilon)+W_{\rm load}+W_{\rm output}$.
\end{proposition}

Uniform primitive moduli are a substantive condition. For example, a recurrence $L_t\le C+\bar\beta F L_{t+1}$ has a horizon-uniform bound when $\bar\beta F<1$ and the terminal modulus is bounded. Without this condition the required covers can themselves grow with the horizon. For undiscounted finite problems, uniform local allowances of order $\varepsilon/(T+1)$ still give a valid construction, but the resolution work generally grows with $T$. Neither statement removes dimensional or precision costs.

The bound gives a constructive response to descendant growth, rather than a reinterpretation of the old recursive timings. The stored-reference construction in Section~\ref{sec:accuracy62} is one concrete backend; the exact neural envelope is another. They must be compared at a common original-optimum tolerance and with their preprocessing included. The retained four-date vector calculation is not relabeled as a successful execution of either alternative.

\subsection{Optimal allocation under a declared work model}
Let $w_t>0$ be the policy-loss weight of a date-$t$ numerical allowance $\eta_t$. Suppose a particular refinement stage has a proved separable work bound $\kappa_t\eta_t^{-p}$, with $p>0$. This model can describe one component of \eqref{eq:work86}; it does not absorb other components whose cost depends jointly on adjacent resolutions.

\begin{theorem}[Weighted accuracy allocation]\label{thm:allocation86}
For a finite collection of stages, define
\[
 S=\sum_t\kappa_t^{1/(p+1)}w_t^{p/(p+1)}.
\]
Among positive allowances with $\sum_tw_t\eta_t\le\varepsilon$, the minimum separable work and its unique allocation are
\begin{equation}\label{eq:allocation86}
 W_* = S^{p+1}\varepsilon^{-p},\qquad
 \eta_t^*=\frac{\varepsilon}{S}
                  \left(\frac{\kappa_t}{w_t}\right)^{1/(p+1)}.
\end{equation}
If an implementation imposes bounds on the admissible allowances, the constrained convex program must be solved with those bounds; the unconstrained formula is not a feasible certificate by itself.
\end{theorem}

For $w_t=\bar\beta^t$ and uniformly bounded $\kappa_t$, $S$ is bounded by a geometric series independently of $T$. This bounds the stated refinement component only. At least the per-date model evaluation, loading and output in \eqref{eq:work86} must still be paid. The exact regression example uses $p=1$, four stages, $w_t=4^{-t}$, $\kappa_t=1$ and $\varepsilon=1/16$. The optimal allowances are $(1/30,1/15,2/15,4/15)$ and their work is $225/4$, compared with $85$ under equal allowances. These are mathematical cost-model values, not measured processor operations.

\subsection{Resource prices and elapsed time}
A machine-independent ledger is a vector rather than a single universal weight. Its coordinates include primitive and endpoint evaluations, descendant expansions, interval arithmetic, comparisons, model evaluation, rational bit operations, bytes loaded, bytes serialized and peak live storage. A price vector $p$ turns that vector $r_j$ into a cost $p^{\mathsf T}r_j$. Complete process time, CPU time and peak resident memory are reported separately. A count does not identify their conversion, and two machines need not induce the same ordering. The inherited reversal between deterministic work and complete elapsed time is consequently retained as evidence about different resource prices, not averaged away.

The current reproducibility experiment records both whole-bin implementations, all acquisition ambiguities, actor queries, exact midpoint checks and full-process timing. It does not claim energy measurements or a calibrated monetary price for a processor operation. Those quantities remain external economic inputs. Section~\ref{sec:decision86} gives conclusions uniformly over their declared ranges, so that a purchaser need not accept an arbitrary scalar operation weight.
